% ===========================================================================
% AITP L4 Test Result Record — Example: QHO Benchmark (outcome: pass)
% ===========================================================================
% This is a filled-in example of the L4 test result record template.
% Topic: Quantum Harmonic Oscillator ground state energy benchmark
% ===========================================================================

\documentclass[11pt,a4paper]{article}

\usepackage[utf8]{inputenc}
\usepackage[T1]{fontenc}
\usepackage{amsmath,amssymb,amsfonts}
\usepackage{physics}
\usepackage{bm}
\usepackage{geometry}
\geometry{margin=2.5cm}
\setlength{\parindent}{0pt}
\setlength{\parskip}{0.5em plus 0.2em minus 0.1em}
\usepackage{booktabs}
\usepackage{longtable}
\usepackage{array}
\usepackage{tabularx}
\usepackage{enumitem}
\usepackage{xcolor}
\usepackage[most]{tcolorbox}
\usepackage{hyperref}
\hypersetup{colorlinks=true,linkcolor=blue!60!black,urlcolor=blue!60!black}

\definecolor{passgreen}{RGB}{34,139,34}
\definecolor{partialorange}{RGB}{204,119,34}
\definecolor{failred}{RGB}{178,34,34}
\definecolor{stuckgray}{RGB}{128,128,128}
\definecolor{infoblue}{RGB}{47,85,151}

\newcommand{\outcomepass}[0]{\textcolor{passgreen}{\textbf{PASS}}}
\newcommand{\outcomepartial}[0]{\textcolor{partialorange}{\textbf{PARTIAL\_PASS}}}
\newcommand{\outcomefail}[0]{\textcolor{failred}{\textbf{FAIL}}}
\newcommand{\checkpass}[0]{\textcolor{passgreen}{$\checkmark$ PASS}}
\newcommand{\checkwarn}[0]{\textcolor{partialorange}{$\sim$ INCONCLUSIVE}}
\newcommand{\checkfail}[0]{\textcolor{failred}{$\times$ FAIL}}

\newtcolorbox{infobox}[1][]{colback=infoblue!4,colframe=infoblue!70,fonttitle=\bfseries,boxrule=0.6pt,arc=2pt,breakable,title={#1}}
\newtcolorbox{passbox}[1][]{colback=passgreen!4,colframe=passgreen!70,fonttitle=\bfseries,boxrule=0.6pt,arc=2pt,breakable,title={#1}}
\newtcolorbox{warnbox}[1][]{colback=partialorange!4,colframe=partialorange!70,fonttitle=\bfseries,boxrule=0.6pt,arc=2pt,breakable,title={#1}}
\newtcolorbox{failbox}[1][]{colback=failred!4,colframe=failred!70,fonttitle=\bfseries,boxrule=0.6pt,arc=2pt,breakable,title={#1}}

\title{\textbf{L4 Validation Test Result Record}\\[0.4em]
  \large QHO Ground State Energy Benchmark}
\author{AITP Protocol v3 — L4 Test Record}
\date{2026-04-20}

\begin{document}
\maketitle

\begin{abstract}
This document records one complete L4 validation run for the quantum harmonic
oscillator ground state energy benchmark. Validation surface: toy\_numeric.
Lane: formal\_theory with numerical cross-check. Outcome: PASS.
\end{abstract}

\tableofcontents
\newpage

% ===========================================================================
\section{Run Identification}

\begin{infobox}[Run Metadata]
\begin{tabularx}{\linewidth}{@{}>{\raggedright\arraybackslash}p{3.5cm} >{\raggedright\arraybackslash}X@{}}
\toprule
\textbf{Field} & \textbf{Value} \\
\midrule
Topic slug       & \texttt{qho-benchmark} \\
Run ID           & \texttt{run-2026-04-20-a} \\
Lane             & \texttt{formal\_theory} \\
Date             & \texttt{2026-04-20T12:00:00Z} \\
L4 cycle         & 1 \\
Validation surface & \texttt{toy\_numeric} \\
\bottomrule
\end{tabularx}
\end{infobox}

% ===========================================================================
\section{Candidate Under Test}

\begin{infobox}[Candidate Summary]
\textbf{Candidate ID:} \texttt{cand-qho-energy-spectrum}

\textbf{Claim:} The energy eigenvalues of the 1D quantum harmonic oscillator
with Hamiltonian $H = p^2/(2m) + \frac{1}{2} m \omega^2 x^2$ are
$E_n = (n + 1/2) \hbar \omega$, for $n = 0, 1, 2, \ldots$

\textbf{Type:} \texttt{result} \quad
\textbf{Regime:} Non-relativistic QM, 1D, harmonic potential only
\end{infobox}

% ===========================================================================
\section{Validation Plan}

\subsection{Mandatory Checks}

\begin{longtable}{>{\raggedright}p{4cm} p{8cm}}
\toprule
\textbf{Check} & \textbf{Plan} \\
\midrule
\endhead
dimensional\_consistency & Verify $[E_n] = [\hbar\omega] = \text{energy}$. Use SymPy dim analysis. \\
symmetry\_compatibility  & Check $[H, \hat{P}] = 0$ where $\hat{P}$ is parity. \\
limiting\_case\_check    & Verify classical HO period $T = 2\pi/\omega$ recovered as $n \to \infty$. \\
correspondence\_check    & Compare $E_n$ against Griffiths QM \S2.3 and known experimental values. \\
\bottomrule
\end{longtable}

\subsection{Pass/Fail Criteria}

\begin{infobox}[Pass Conditions]
All four mandatory checks must return PASS. Dimensional check must confirm
$ML^2T^{-2}$ on both sides. Correspondence check must match within 1\%.
\end{infobox}

\begin{failbox}[Failure Signals]
Dimensional mismatch, non-commutation with parity, classical limit not
recovered, or $>5\%$ deviation from known $E_n$ values.
\end{failbox}

% ===========================================================================
\section{Physics Check Results}

\subsection{Check Result Summary}

\begin{longtable}{>{\raggedright}p{4cm} p{2.5cm} p{6cm}}
\toprule
\textbf{Check} & \textbf{Verdict} & \textbf{Detail} \\
\midrule
\endhead
dimensional\_consistency & \checkpass & $[H] = [\hbar\omega] = ML^2T^{-2}$ both sides \\
symmetry\_compatibility  & \checkpass & $[H, \hat{P}] = 0$ confirmed \\
limiting\_case\_check    & \checkpass & $T = 2\pi/\omega$ recovered as $n \to \infty$ \\
correspondence\_check    & \checkpass & $E_n = (n+1/2)\hbar\omega$ matches Griffiths \S2.3 \\
\bottomrule
\end{longtable}

\subsection{SymPy Verification Evidence}

\textbf{Tool:} \texttt{aitp\_verify\_dimensions}

\begin{verbatim}
{
  "pass": true,
  "lhs_dimension": [1, 2, -2, 0, 0],
  "rhs_dimension": [1, 2, -2, 0, 0],
  "expression_lhs": "H",
  "expression_rhs": "hbar * omega * (n + 1/2)"
}
\end{verbatim}

% ===========================================================================
\section{Numerical Results}

For the ground state ($n=0$):
\begin{equation}
E_0 = \frac{1}{2}\hbar\omega
\end{equation}

\textbf{Benchmark comparison:}
\begin{itemize}
  \item \textbf{Analytic:} $E_0 = 0.5\,\hbar\omega$
  \item \textbf{Numerical (exact diagonalization, $N_\text{basis}=100$):} $E_0 = 0.5000000\,\hbar\omega$
  \item \textbf{Relative error:} $< 10^{-12}$ (machine precision)
\end{itemize}

No external compute was required for this validation since the result is
analytic. For numerical benchmarks, HPC job details would appear here.

% ===========================================================================
\section{Execution Provenance}

\begin{infobox}[Execution Environment]
\begin{tabularx}{\linewidth}{@{}>{\raggedright\arraybackslash}p{3.5cm} >{\raggedright\arraybackslash}X@{}}
\toprule
\textbf{Field} & \textbf{Value} \\
\midrule
Host            & \texttt{local (macOS ARM64)} \\
Job ID          & \texttt{N/A (analytic validation)} \\
Walltime used   & N/A \\
Code version    & \texttt{AITP v3, SymPy 1.12} \\
Compiler        & N/A \\
MPI processes   & N/A \\
\bottomrule
\end{tabularx}
\end{infobox}

\subsection{Data Provenance}

\begin{longtable}{>{\raggedright}p{2.5cm} p{3cm} p{2cm} p{4.5cm}}
\toprule
\textbf{Data Point} & \textbf{Script} & \textbf{Executed At} & \textbf{Method} \\
\midrule
\endhead
dim\_analysis & aitp\_verify\_dimensions & 2026-04-20T11:55Z & SymPy dimensional analysis \\
commutation   & aitp\_verify\_algebra    & 2026-04-20T11:56Z & SymPy commutator evaluation \\
limiting\_case & analytic\_limit.py       & 2026-04-20T11:57Z & Analytic limit $n\to\infty$ \\
\bottomrule
\end{longtable}

% ===========================================================================
\section{Devil's Advocate}

\begin{warnbox}[Adversarial Review]
The derivation assumes the potential is exactly harmonic. Anharmonic
corrections ($x^3$, $x^4$ terms) would shift the spectrum. Regime boundary:
coupling must be weak enough that perturbation theory converges.

Additionally:
\begin{itemize}
  \item The result is for a single particle in 1D. Extension to $N$ interacting
    particles in 3D requires independent verification.
  \item The classical limit check only verifies the period, not the full
    classical phase space structure.
  \item The correspondence check relies on a single textbook source. Cross-check
    with an independent method (path integral) would strengthen confidence.
\end{itemize}
\end{warnbox}

% ===========================================================================
\section{Comparison to Known Results}

\begin{longtable}{>{\raggedright}p{3cm} p{3cm} p{3cm} p{3.5cm}}
\toprule
\textbf{Source} & \textbf{Result} & \textbf{Our Result} & \textbf{Agreement} \\
\midrule
\endhead
Griffiths QM \S2.3 & $E_n = (n+1/2)\hbar\omega$ & $E_n = (n+1/2)\hbar\omega$ & Exact \\
Sakurai \S2.3      & $E_n = (n+1/2)\hbar\omega$ & $E_n = (n+1/2)\hbar\omega$ & Exact \\
Cohen-Tannoudji    & $E_n = (n+1/2)\hbar\omega$ & $E_n = (n+1/2)\hbar\omega$ & Exact \\
\bottomrule
\end{longtable}

% ===========================================================================
\section{Final Outcome}

\begin{passbox}[L4 Outcome]
\textbf{\Large \outcomepass}

All four mandatory physics checks passed. The candidate survives adversarial
review. No contradictions with known results. No capability gaps.
\end{passbox}

\subsection{Promotion Decision}

\begin{itemize}
  \item \textbf{Verdict:} \texttt{promote}
  \item \textbf{Trust basis:} \texttt{validated}
  \item \textbf{Trust scope:} \texttt{bounded\_reusable}
  \item \textbf{Promoted to L2 node:} \texttt{qho-energy-spectrum}
  \item \textbf{Decider:} AITP L4 adjudication + operator approval
  \item \textbf{Decision timestamp:} 2026-04-20T12:10:00Z
\end{itemize}

% ===========================================================================
\section{Backed-Up Data}

\begin{infobox}[Backup Manifest Summary]
\begin{tabularx}{\linewidth}{@{}>{\raggedright\arraybackslash}p{3.5cm} >{\raggedright\arraybackslash}X@{}}
\toprule
\textbf{Class} & \textbf{Files} \\
\midrule
Test definition  & \texttt{validation\_contract.md}, \texttt{execution-tasks/dim\_check.json} \\
Key inputs       & \texttt{convention\_snapshot.md}, \texttt{derivation\_anchor\_map.md} \\
Final outputs    & \texttt{dim\_check\_result.json}, \texttt{algebra\_verify\_result.json} \\
Key logs         & \texttt{validation.log} \\
Figures          & (none — analytic validation) \\
\bottomrule
\end{tabularx}
\end{infobox}

\noindent\textbf{Backup location:} \texttt{backups/l4/qho-benchmark/run-2026-04-20-a/}

\noindent\textbf{Original size:} 0.5 MB

\noindent\textbf{Backup size:} 0.05 MB

\noindent\textbf{Removed classes:} scratch notes, intermediate SymPy dumps, raw exploration files

% ===========================================================================
\section{Open Issues \& Follow-Up}

\subsection{Known Limitations}
\begin{itemize}
  \item Harmonic approximation only. Anharmonic corrections not quantified.
  \item Single-particle result. Many-body extension not validated.
  \item 1D only. Higher-dimensional harmonic oscillator not checked.
\end{itemize}

\subsection{Follow-Up Tasks}
\begin{itemize}
  \item Estimate first-order anharmonic correction $\Delta E_0^{(1)}$ for
    $V(x) = \frac{1}{2}m\omega^2 x^2 + \lambda x^4$ (route to L3).
  \item Extend validation to 3D isotropic harmonic oscillator.
  \item Cross-check with path integral derivation (independent method).
\end{itemize}

\subsection{Gap Classification}
No gaps. All mandatory checks closed.

% ===========================================================================
\end{document}
